% Concise front section for the 15-page SpaceCHI review manuscript.
% The named manuscript retains the expanded exposition in sections 01--03.
\section{Introduction}

Cislunar space is the Earth--Moon region. A trajectory-correction maneuver
changes a spacecraft's velocity to keep its path within mission constraints;
its effort is expressed as delta-v (m/s), which is related to but is not itself
propellant mass. For a crew-carrying mission, correction planning couples burn
timing and delta-v with navigation uncertainty, execution error, crew
timelines, entry conditions, and reserves. Public Artemis trajectory-design
research consequently evaluates robust correction cost under modeled
uncertainty and constraints \cite{woffinden2023}. A learned surrogate could
screen candidate scenarios faster than repeated high-fidelity simulation, but
that possibility is a motivation for testing rather than evidence of mission
utility.

This paper treats such screening as human-centered decision support. A model
score is not a decision: a qualified reviewer must be able to inspect the data,
physical assumptions, uncertainty boundary, comparators, and failure cases
before authorizing further study. Crew-related requirements therefore remain
separate from predictive performance, consistent with NASA model-credibility
and human-system standards \cite{nasa7009,nasa3001}. No user study, interface
evaluation, or measurement of trust, workload, or operator performance was
conducted.

Quantum machine learning (QML) places quantum-state representations or
circuits inside a learning method. A quantum kernel encodes each classical
input as a quantum state and uses state overlap as a similarity measure
\cite{schuld2019,havlicek2019}. A distinct feature space, however, does not
establish useful prediction or quantum advantage. Data loading, encoding
geometry, classical approximability, split leakage, weak comparators, and
post-outcome selection can all produce misleading conclusions
\cite{huang2021,sweke2025,bowles2024,schnabel2024}. These are evidence-design
problems before they are quantum-computing problems.

\subsection{Research question and contributions}

The central question is: \emph{To what extent can simulated quantum-kernel
representations add predictive value to a human-centered cislunar
trajectory-correction benchmark when compared fairly with physics-based
controls?} The frozen program tested a primary fidelity kernel (Q01), a
projected follow-up (Q01b), a feasibility-only kernel (FQK), and bounded
successors. It also asked what could be concluded while calibration,
final-test, mission-loop, and hardware evidence remained locked. The tested
models answered the performance questions negatively; the locked evidence
defines the scope of that answer.

The paper contributes five connected artifacts:
\begin{enumerate}
\item a source-bound benchmark with grouped folds and transformations fitted
only on each training fold;
\item a simulator-credibility chain that retains its source-ineligible
exception rather than converting missing evidence into a pass or failure;
\item matched physics-residual, random-feature, compressed-input, and
classical-kernel controls;
\item a complete ledger of the primary, exploratory, and bounded-successor
outcomes, including valid negative results; and
\item an evidence package that keeps model results, human authorization, and
operational claims distinct.
\end{enumerate}
Together these artifacts make a rejected advancement hypothesis inspectable
instead of treating a negative result as failed record keeping.

\section{Scientific Scope and Claim Boundary}
\label{sec:scope}

This is an offline, public-data, development-only benchmark, not mission-design
authority, flight certification, or quantum-hardware assessment. Repository
identifiers are audit keys: a Gate is an authorization checkpoint; D and P
identify decisions and prospective protocols; C, Q, and A identify classical,
QML, and analysis controls; and RFIG identifies figure provenance. C06 is the
physics-residual control, Q01 the primary kernel, and D046 a bounded successor.
A fold keeps related cases together, a seed fixes a random start, and a rung is
a 128--1,024-row sample stage. Development rows may select models; calibration
would adjust a frozen output and final-test rows would estimate performance
after all choices were fixed. Neither locked split was read.

The intended role is offline screening information for qualified human review,
with physical analysis as the fallback. Gate 3 supports only eligible simulator
comparisons; Gate 5 supports a frozen development decision; later experiments
cannot revise Gate 5 or authorize Gate 6. The recall-first classical analysis
is a future design lesson, not a safety filter. Mission-loop use,
quantum-hardware or GPU execution, operational thresholds, and trained-model
release remain outside scope. The paper claims no quantum advantage,
propellant saving, flight readiness, NASA approval, hardware behavior,
final-test result, or QML invention.

Figure \ref{fig:timeline} records authorization state rather than model
performance.

\begin{figure}[!t]
\centering
\includegraphics[width=0.76\textwidth]{figures/gate_decision_timeline.png}
\caption{Governance through Gate 5 (RFIG-001). Gates 0--5 were closed by
human decisions; Gate 5 rejected a benchmark-specific advancement hypothesis,
and Gate 6 was not opened.}
\label{fig:timeline}
\end{figure}

\section{Related Work and Positioning}

The benchmark is positioned against public Artemis correction-burn work, not
as a replacement for mission-owned analysis. Woffinden et al. combine
uncertainty, timing, entry, and crew-related constraints \cite{woffinden2023}.
Classical learning studies of low-thrust guidance and cislunar transfers offer
precedents for surrogate modeling and feasibility handling, but differ in
dynamics, propulsion, objectives, and validation structure
\cite{zavoli2020,sullivan2021,izzo2021,steffen2025,xie2022,zhang2025}.

Quantum feature-space methods provide alternative similarities, but their
value depends on whether that geometry reflects the target and remains useful
against strong classical models \cite{schuld2019,havlicek2019,huang2021}.
Encoding scale and bandwidth alter inductive bias \cite{shaydulin2022}; random
features provide a relevant distance-kernel control \cite{rahimi2007}; and
fair comparisons require explicit splits, seeds, capacities, and selection
rules \cite{bowles2024,schnabel2024,pernot2020}. Data re-uploading,
generalization, and initialization studies informed the frozen encoding,
observable, and seed choices \cite{perez2020,caro2021,grant2019}. Adjacent
quantum-agent and trajectory-optimization work informed controls and scope but
was not reproduced here \cite{skolik2022,degrossi2025}.

The repository records a bounded set of 23 accepted literature sources plus
the random-Fourier source behind control A01. It distinguishes methodological
influence from methods actually tested. A later archival search did not change
the protocol or outcomes; this paper therefore makes a benchmark claim, not an
exhaustive-review claim.
